
\begin{figure}[H]
    \centering
    \begin{tabular}{c}
        % Corrupted Input (z_0)
        \fbox{\begin{minipage}{0.9\textwidth}
            \tiny <|endoftext|>When her child was planning to travel to Mexico, she thought it could be easy for her to find money to pay the trip. But that didn't happen, either. The girl arrived on her doorstep the day before after they were scheduled to return to home and - despite her own passport blocked. It wasn't easy. Even her mother reported a bad visit. A gate to the airport- the gate and flight instructions - were locked at nine a.. "I had no way to the gate," she told The Local. "I feel like I should keep my passport and their fees." At first the police arrived, a British man told an e-mail that parents should only have money for themselves to be fine, but could never explain the situation. Later, she told her daughter that her passport had been detained because she wanted to visit her family at a hotel. She felt embarrassed and confused. Nearly two months later, she said, she lost $35 ($40) on rent and imagined that the traffic-police and customs agents in Bangkok would end up delaying flights and forcing her to stay home. She was worried that her father would refuse allowing her daughter to spend a few days in the country. Meanwhile, the police were sent to search. "It wasn't easy for them when a child feels like home for the first time," said Mahavram Kaas, a spokesman for the Ministry of Foreign Affairs. He's referring to a tour arranged by the French and French foreign ministry, known as Courage in the Child. That tour cost the region 2 million worth of tickets, and cost the Calais family about \$25 million in lodging expenses, according to a statement by both the ministry's behalf (he told the Portuguese police agents) and London's Embassy in London (he told the French ambassador they provided a payment for 70,000, which would be used to pay the travel costs for their visits). In 2011, a family from Calais had moved to the UK aged 15. Their sister left to remain in France at 5 years old. Her brother tried to answer that question. He explained it to reporters at his service station at Calais airport. "You take the morning. It's named after you and your little girls," he said. The last mother was having a 19-month stay with her daughter that night, the police said. The first time she was back her husband took their daughter on a boat to the UK, said the mayor of the British Transport Agency. That meant she had plenty of cash-to-go and no money to borrow when she opened an account at Kathmandu Airport a few hours after booking her flight. She panicked. She called the was a Daley's Nessie (small cash register), saying she was getting better. Her doctors visited her when she returned and her boyfriend quit his job for four months after the visit, she said. "How do you feel like you are safe?" A text from a friend left her to the police. "She says I must go get my wallet," said Ajaz. Soon after, she booked a plane ticket to Paris and took a metro train to Calais. One night, a French policeman would knock on the door of a local council building, open the mailman and the phone and tell her she knew that she could not leave at least one week without food. The four months her daughter spent in the UK was exhausting and hard, and it reached the stage where she realized she could barely stay in Britain. "Now no choice but to go home. Then we regret having a daughter," he said. He thought for a minute. "You've broken your heart." "Today, my daughter and my boyfriend decided to stay in this country for over two months," he wrote in an email with his daughter in his hand. "All our flights cancelled and no security. Shame." Caines' family were also put on leave. The French police paid for her car after she rented it, and her female officer used it for the opening ceremony of her press conference in Thailand. The police are still arranging for her family to have their official visit. Although her son is back at work now and his old job, her daughter needs to stay in a hospital in Algiers to continue her education. But, finally, her parents will be making their girl home. The daughter was 18 when they opened her case. She was born two months ago. She doesn't talk about it because it feels like she was still a child, living in Thailand with a small child. Her mother, Anzsa Gurdon, came to England as a three-year-old after her brother, Ehab Rahman, was working as a British worker in Calais while living in London and studying abroad in Dubai. As a mother earns 2,000 pounds a month, they receive a well-paid living in secure accommodation, some even with public transport buses. If they make money, their child stays in the UK, they can set up companies with kids to take care of their children. Other countries sometimes also give birth to parents forced to provide child care. The parents are often refugees from their home countries, they're left without family, and have forced to leave families. As one former refugee fled Syria, his family was in detention, because when and if their child had arrived, they would be living somewhere. The detention centers in Western Europe often have a higher rate for asylum seekers. Often there are higher "safe houses" for young people, and then the people in the center get older when their child comes to stay, but the only family that is a year older is not allowed to have children, and usually only if they stay six months. Children are also detained and are asked to show identification. Because in most cases the refugees ask only about their identity, they don't have access to their own documents, and have no other documentation. An activist working for Ireland says he's against offshore processing. He thinks the charity problem here is like diseases which don't sufficiently seek out international funding. I think too many countries want to employ "humanitarian" children. Now the job Before the refugee crisis in Calais, 1,823 children were living in the UK, the UNHCR website shows. A good chunk of those children had landed in refugee camps in Africa, where mostly African migrants were sending children from Syria and Libya to their camps, but those numbers didn't fare so well. "At the moment when I met the French, it was horrific, they wanted me to put my children in a van. But I was only kidding. It's something called 10-year vans," Aakaz said during an interview. "I want to keep my children for 10 years. That's something. It's like Christmas. The dream of 10 years. . . . For me, the idea that this is a good opportunity, here's a chance," is that she can sleep with her children. At this point, it's much more than just about "alternatives". They now have to decide, at some point, whether they want to take the chance. Is it a big deal or not? Only in Calais She explained what's agreed to so the children can go home and can have a better future. "I'm very determined. My children want to go home but it's my life's personal decision," she said. "Five months. I want them to be home, 5 months. If they're not getting jobs properly, I also want to stay home. But I feel good about what I've got. She is from a poor country. I don't owe anything to anyone. But I have to work for them. I feel like I can just go to the road and provide accommodation for my children and their children." But they all don't work well because their families have her as a head. "They want me to have a job in England, but I feel like it's my home, and I'm not scared of work," she said. "And I feel that the opposite would be possible. I think in the future that I can have a job or two there." A Second World Friday event will be held outside Calais on Saturday, donating 100 euros to the coming week in the money brought up to them by UNHCR through the King Wahab Samba Global Fund. Friends and family expressed their "weakness" like many survivors in many countries."We just don't want to accept what has to happen. We want to put the people back there as soon as possible," said one man. Her brother, who is the son of a long term migrant, said: "The story of the refugee is not a mother's story. The story of the refugee is children's story." "In Calais it's too young for these kids. They play outside or work outside, they just eat, right? I don't think much has changed. This child with all her food and sleep, she's too young for life without any protection. We don't need any protection at all. We need anything that would be safer."<|endoftext|>
    \end{minipage}} \\
    \\
    \end{tabular}
    \caption{블록 크기 $L'=16$의 \algo{}에서 길이 $L=2031$의 샘플($T=5$K 디퓨전 단계, 문맥 길이 $L=1024$로 학습). GPT2-Large 하에서 이 샘플의 생성 perplexity는 24.3, 엔트로피는 5.5이다. (생성된 영어 텍스트는 평가 대상이므로 원문 그대로 둔다.)}
    \label{sar-samples}
\end{figure}


